\chapter{पैसा और माप}\label{ch:g2:measure}

मापपट्टियाँ, रसोई का तराज़ू, जग, सिक्के, घड़ियाँ और कैलेंडर:
यह अध्याय रोज़मर्रा की संख्याओं का है --- कितना लंबा, कितना
भारी, कितना समाता है, कितने का है, और कितने बजे
हैं।

\section{लंबाई, द्रव्यमान, धारिता}

\begin{definition}[रोज़ के मात्रक]\label{def:g2:measure:units}
\begin{itemize}
  \item \emph{लंबाई}: छोटी चीज़ों के लिए सेंटीमीटर (cm), बड़ी
        चीज़ों के लिए मीटर (m), और $1$ m $= 100$ cm;
  \item \emph{द्रव्यमान}\index{द्रव्यमान}: हल्की चीज़ों के लिए ग्राम (g), भारी
        के लिए किलोग्राम (kg), और $1$ kg $= 1000$ g;
  \item \emph{धारिता}\index{धारिता} (कोई पात्र कितना समाता है): लीटर (L)।
\end{itemize}
इससे छोटे मात्रक और उनके पक्के नियम \cref{ch:g3:measure} में आएँगे।
\end{definition}

\begin{example}[कौन-सा मात्रक?]\label{ex:g2:measure:choose}
एक रंगीन पेंसिल: लगभग $9$ cm। एक दरवाज़ा: लगभग $2$ m। एक सेब:
लगभग $150$ g। आटे का बड़ा थैला: $1$ kg। दूध की बोतल: $1$ L।
मापने से पहले मात्रक का अंदाज़ा लगा लो --- इससे ``दरवाज़ा $2$
cm का है'' जैसे बेतुके उत्तर नहीं आते।
\end{example}

\section{पैसा}

\begin{method}[रक़म बनाना]\label{met:g2:measure:money}
$1$, $2$, $5$, $10$, $20$ के सिक्कों और नोटों से:
\begin{enumerate}
  \item जो सबसे बड़ा समा जाए उससे शुरू करो, फिर बाक़ी भरो;
  \item $27$ बनाने के लिए: एक $20$, एक $5$ और एक $2$;
  \item अक्सर कई तरीक़े होते हैं --- $27$, $10 + 10 + 5 +
        2$ भी है।
\end{enumerate}
\end{method}

\begin{omfigure}
\begin{tikzpicture}[scale=0.9]
  \draw[thick, fill=omProp!20] (0,0) rectangle (1.3,0.7);
  \node[font=\small] at (0.65,0.35) {$20$};
  \node[font=\large] at (1.9,0.35) {$+$};
  \draw[thick, fill=omMeth!20] (2.85,0.35) circle (0.35);
  \node[font=\small] at (2.85,0.35) {$5$};
  \node[font=\large] at (3.8,0.35) {$+$};
  \draw[thick, fill=omDef!20] (4.7,0.35) circle (0.3);
  \node[font=\small] at (4.7,0.35) {$2$};
  \node[font=\large] at (5.6,0.35) {$=$};
  \node[font=\normalsize] at (6.5,0.35) {$27$};
\end{tikzpicture}

{\small एक नोट और दो सिक्कों से $27$ --- और दूसरे तरीक़े भी हैं,
जैसे $10 + 10 + 5 + 2$।}
\end{omfigure}

\begin{example}[बाक़ी पैसे]\label{ex:g2:measure:change}
एक किताब $14$ की है; तुम $20$ देते हो। बाक़ी पैसे $20 - 14 =
6$ हैं, जो आगे गिनकर मिलते हैं: $14 \to 20$ पर $6$
(\cref{met:g2:subtraction:countup})।
\end{example}

\section{समय}

\begin{method}[आधा और चौथाई घंटा]\label{met:g2:measure:clock}
लंबी सुई बताती है कि घंटे के बाद कितने मिनट हुए:
\begin{enumerate}
  \item सीधे ऊपर ($12$): पूरा घंटा --- $3{:}00$;
  \item दाईं ओर चौथाई घुमाव ($3$): सवा --- $3{:}15$, ``सवा तीन'';
  \item सीधे नीचे ($6$): साढ़े --- $3{:}30$, ``साढ़े तीन'';
  \item तीन-चौथाई घुमाव ($9$): अगले घंटे में चौथाई बाक़ी ---
        $3{:}45$, ``पौने चार''।
\end{enumerate}
\end{method}

\begin{omfigure}
\begin{tikzpicture}[scale=0.8]
  \foreach \x/\hh/\mm/\lb in {0/3/3/{3:15}, 4/3/6/{3:30},
                                8/3/9/{3:45}} {
    \begin{scope}[xshift=\x cm]
    \draw[thick] (0,0) circle (1.2);
    \foreach \h in {1,...,12} {
      \node[font=\tiny] at ({90-\h*30}:0.98) {\h};
    }
    \draw[very thick, omDef] (0,0) -- ({90-(\hh+\mm/12)*30}:0.55);
    \draw[thick, omThm] (0,0) -- ({90-\mm*30}:0.95);
    \node[below, font=\small] at (0,-1.45) {\lb};
    \end{scope}
  }
\end{tikzpicture}

{\small सवा, साढ़े, पौने: लंबी लाल सुई $3$, $6$ और
$9$ पर।}
\end{omfigure}

\begin{example}[कैलेंडर]\label{ex:g2:measure:calendar}
एक साल में $12$ महीने होते हैं, क्रम से: जनवरी, फ़रवरी, मार्च,
अप्रैल, मई, जून, जुलाई, अगस्त, सितंबर, अक्टूबर, नवंबर, दिसंबर।
ज़्यादातर महीनों में $30$ या $31$ दिन होते हैं; फ़रवरी में $28$
(या $29$)। तारीख़ दिन और महीना बताती है: ``$5$ मार्च''।
\end{example}

\section{अभ्यास}

\begin{exercise}[$\star$]\label{exo:g2:measure:1}
मात्रक चुनो (cm, m, g, kg या L): एक पेंसिल; कक्षा की
लंबाई; एक पंख; एक बस्ता; पानी की बोतल।
\end{exercise}

\begin{exercise}[$\star$]\label{exo:g2:measure:2}
बदलो: $2$ m को cm में; $300$ cm को m में; $1$ kg को g में।
\end{exercise}

\begin{exercise}[$\star$]\label{exo:g2:measure:3}
सिक्कों और नोटों से ये रक़में बनाओ: $16$; \quad $23$; \quad
$35$। हर एक का एक तरीक़ा बताओ।
\end{exercise}

\begin{exercise}[$\star$]\label{exo:g2:measure:4}
बाक़ी पैसे निकालो: $6$ के खिलौने के लिए $10$ दो; $17$ की किताब
के लिए $20$; $32$ के खेल के लिए $50$।
\end{exercise}

\begin{exercise}[$\star$]\label{exo:g2:measure:5}
समय पढ़ो: लंबी सुई $6$ पर, छोटी सुई $4$ और $5$ के बीच;
लंबी सुई $3$ पर, छोटी सुई $8$ से थोड़ा आगे।
\end{exercise}

\begin{exercise}[$\star$]\label{exo:g2:measure:6}
साढ़े दो बजे दिखाती घड़ी बनाओ। हर सुई कहाँ है?
\end{exercise}

\begin{exercise}[$\star$]\label{exo:g2:measure:7}
महीना बताओ: मई के बाद; जनवरी से पहले; अपने जन्मदिन का
महीना।
\end{exercise}

\begin{exercise}[$\star$]\label{exo:g2:measure:8}
लेआ $3$ का पेय और $4$ का केक ख़रीदती है। वह $10$ का नोट देती
है। बाक़ी कितने मिलेंगे?
\end{exercise}

\begin{exercise}[$\star$]\label{exo:g2:measure:9}
कौन भारी है: $1$ kg चीनी का थैला या $800$ g सेब? कितने ग्राम
से?
\end{exercise}

\begin{exercise}[$\star\star$]\label{exo:g2:measure:10}
ठीक चार सिक्कों या नोटों से ($1$, $2$, $5$, $10$, $20$ में से
चुनकर) $40$ बनाओ। दो अलग-अलग तरीक़े ढूँढ़ो।
\end{exercise}
